\documentclass[conference]{IEEEtran}
\IEEEoverridecommandlockouts

% Paquetes esenciales para compilación en Overleaf
\usepackage[spanish,es-tabla]{babel}
\usepackage[utf8]{inputenc}
\usepackage[T1]{fontenc}
\usepackage{cite}
\usepackage{amsmath,amssymb,amsfonts}
\usepackage{algorithmic}
\usepackage{graphicx}
\usepackage{textcomp}
\usepackage{xcolor}
\usepackage{booktabs}
\usepackage{tabularx}
\usepackage{url}
\usepackage{hyperref}

\hypersetup{
    colorlinks=true,
    linkcolor=blue,
    filecolor=magenta,      
    urlcolor=cyan,
    citecolor=blue
}

\def\BibTeX{{\rm B\kern-.05em{\sc i\kern-.025em b}\kern-.08em
    T\kern-.1667em\lower.7ex\hbox{E}\kern-.125emX}}

\begin{document}

\title{Diseño e Implementación de Infraestructura de Red Local Autónoma para Atención Móvil Comunitaria en Entornos Rurales\\
\large Plan de Proyecto -- Fase 1}

\author{\IEEEauthorblockN{José David Rodríguez Pinilla}
\IEEEauthorblockA{\textit{Código: A00407151} \\
\textit{Profesor: Andrés Navarro Cadavid, Ph.D.} \\
\textit{Infraestructura de Redes 2} \\
\textit{Departamento de Tecnologías de Información y Comunicaciones} \\
\textit{Universidad Icesi}\\
Cali, Colombia \\
jose.rodriguez31@correo.icesi.edu.co}
}

\maketitle

\begin{abstract}
Las zonas rurales y de difícil acceso a menudo carecen de conectividad confiable a Internet o sufren interrupciones eléctricas recurrentes. El presente trabajo detalla el plan de ingeniería para el despliegue de una infraestructura de red local de alta resiliencia diseñada para operar en un vehículo o unidad móvil de atención comunitaria. La solución integra un nodo central de procesamiento de bajo consumo (Mini PC / Raspberry Pi 5), segmentación de red en Capa 2 y Capa 3 mediante portal cautivo para tres grupos de interés (Comunidad, Personal de Salud y Educadores), una arquitectura de almacenamiento masivo local (2--4 TB) y un esquema avanzado de proxy/caché con Squid con políticas de refresco transparentes y sincronización oportunista. Se presenta la topología de red detallada, el direccionamiento IP, el análisis comparativo de hardware, la estrategia de caché híbrida y la ruta crítica de implementación con su respectiva matriz de mitigación de riesgos.
\end{abstract}

\begin{IEEEkeywords}
Redes rurales, Portal Cautivo, Squid Proxy, Caché local, Conmutación Offline/Online, Redes comunitarias, L2/L3.
\end{IEEEkeywords}

\section{Introducción}
El acceso a la información y a los servicios digitales básicos en áreas rurales dispersas enfrenta barreras severas: ausencia de cobertura móvil de alta velocidad, altos costos de transmisión satelital e inestabilidad del fluido eléctrico. En este contexto, una unidad móvil de atención comunitaria requiere de una infraestructura de red independiente que permita ofrecer servicios asistenciales, educativos y de ocio tanto en presencia de un enlace satelital/móvil intermitente (modo \textit{online}) como en aislamiento total (modo \textit{offline}).

El presente documento corresponde a la entrega de la \textbf{Fase 1: Plan de Proyecto}, cuyo propósito es estructurar formalmente la arquitectura técnica de red, la estrategia de almacenamiento/caché y la planificación operativa con el nivel de detalle exigido en la categoría más alta de evaluación (\textit{Accomplished}).

\section{Arquitectura de Red (Capa 2 y Capa 3)}

\subsection{Topología Física y Selección de Hardware}
La red se concibe como una topología en estrella centrada en un dispositivo servidor principal y un punto de acceso (Access Point - AP) inalámbrico de doble banda con capacidad de múltiples SSIDs y etiquetado VLAN (802.1Q). 

En cuanto a la unidad central de cómputo, se contemplan dos alternativas viables:
\begin{enumerate}
    \item \textbf{Raspberry Pi 5 (8 GB RAM):} Arquitectura ARM Cortex-A76 quad-core @ 2.4 GHz, puerto Gigabit Ethernet nativo, interfaz PCIe 2.0 (expandible a almacenamiento NVMe rápido) y consumo energético de 5--12W.
    \item \textbf{Mini PC (Intel Core i5 x86\_64, 16 GB RAM):} Mayor ancho de banda de E/S, doble interfaz Gigabit Ethernet nativa, soporte de virtualización por hardware completo (KVM) y mayor rendimiento para descompresión e indexación de caché Squid.
\end{enumerate}

Para fines de diseño y máxima estabilidad, la arquitectura se parametriza de forma modular para ejecutarse en cualquiera de las dos plataformas mediante contenedores Docker.

\subsection{Diseño en Capa 2 (Segmentación L2 y VLANs)}
Para garantizar aislamiento de tráfico y priorización de ancho de banda, se definen 3 Redes de Área Local Virtual (VLANs) mapeadas a tres SSIDs distintos en el AP:

\begin{itemize}
    \item \textbf{VLAN 10 -- Personal de Salud (SSID: Salud-Atencion):} Tráfico prioritario. Acceso a historias clínicas locales, sincronización de expedientes y conexión segura punto a punto.
    \item \textbf{VLAN 20 -- Educadores y Administración (SSID: Educacion-Rural):} Tráfico educativo. Acceso a plataformas de aprendizaje (Moodle offline, repositorios de libros y videos).
    \item \textbf{VLAN 30 -- Comunidad General (SSID: Comunidad-Libre):} Tráfico restringido por cuotas de tiempo y ancho de banda. Acceso a través de Portal Cautivo para navegación local y caché.
    \item \textbf{VLAN 99 -- Gestión de Infraestructura (Cableada):} Tráfico administrativo exclusivo para configuración local (SSH, consola Nginx/Squid, métricas).
\end{itemize}

\subsection{Diseño en Capa 3 (Direccionamiento IP y Enrutamiento)}
La red adopta el bloque privado \texttt{10.10.0.0/16}, subdividido con VLSM (Variable Length Subnet Masking) para asegurar escalabilidad y separación estricta:

\begin{table}[htbp]
\caption{Esquema de Direccionamiento IP y Segmentación L3}
\begin{center}
\begin{tabularx}{\columnwidth}{l l l X}
\toprule
\textbf{VLAN / Segmento} & \textbf{Subred CIDR} & \textbf{Gateway} & \textbf{Rango DHCP / Capacidad} \\
\midrule
VLAN 10: Salud & 10.10.10.0/24 & 10.10.10.1 & 10.10.10.50 -- 10.10.10.200 (150) \\
VLAN 20: Educación & 10.10.20.0/24 & 10.10.20.1 & 10.10.20.50 -- 10.10.20.200 (150) \\
VLAN 30: Comunidad & 10.10.30.0/23 & 10.10.30.1 & 10.10.30.50 -- 10.10.31.250 (450) \\
VLAN 99: Gestión & 10.10.99.0/28 & 10.10.99.1 & 10.10.99.5 -- 10.10.99.14 (10) \\
WAN / Uplink ISP & DHCP Ext. & Dinámico & Según módem satelital o LTE \\
\bottomrule
\end{tabularx}
\end{center}
\label{tab:ip_addressing}
\end{table}

\subsection{Servicios de Resolución y Enrutamiento (DNS, DHCP y Portal Cautivo)}
El nodo central ejecuta \texttt{dnsmasq} para proveer:
\begin{itemize}
    \item \textbf{DHCP con opciones personalizadas:} Asignación de IPs con lease times diferenciados (2 horas para comunidad, 12 horas para salud).
    \item \textbf{DNS Split-Horizon:} En modo offline, todas las peticiones a dominios de servicios internos (\texttt{portal.local}, \texttt{salud.local}, \texttt{biblioteca.local}) son resueltas a las direcciones virtuales del host.
    \item \textbf{Portal Cautivo (openNDS):} Intercepta el puerto 80/443 de los usuarios de la VLAN 30 mediante reglas \texttt{iptables}, presentando una interfaz de autenticación simple o aceptación de términos antes de autorizar la navegación hacia la caché local o Internet.
\end{itemize}

\section{Estrategia de Caché y Almacenamiento Masivo}

\subsection{Jerarquía de Almacenamiento (2--4 TB)}
El almacenamiento físico consta de un disco externo (SSD o HDD de alta durabilidad mecánica con sistema de archivos \texttt{ext4} y opciones de montaje optimizadas como \texttt{noatime,data=writeback}). La estructura jerárquica se divide de la siguiente manera:
\begin{enumerate}
    \item \textbf{Nivel 1 (Hot Data / RAM \& SSD rápido - 50 GB):} Caché en memoria para objetos ultrarrápidos Squid y base de datos DNS/DHCP.
    \item \textbf{Nivel 2 (Warm Data / Squid Cache Dir - 500 GB a 1 TB):} Almacén estructurado en \texttt{aufs/rock} para páginas web cacheadas, bibliotecas JavaScript/CSS, imágenes y paquetes de software.
    \item \textbf{Nivel 3 (Cold Data / Repositorios Multimedia y Archivos - 2 a 3 TB):} Repositorio estático de enciclopedias offline (Kiwix / ZIM files), videos educativos (Jellyfin), manuales de salud del Ministerio/OMS en formato PDF y réplicas estáticas de sitios web.
\end{enumerate}

\subsection{Políticas de Refresco y Gestión de Squid Proxy}
Squid operará en modo híbrido:
\begin{itemize}
    \item \textbf{Caché Transparente HTTP:} Redirección forzada mediante \texttt{iptables PREROUTING} del puerto 80 hacia el puerto 3128 de Squid.
    \item \textbf{Tráfico HTTPS:} En entornos educativos/offline se implementa un repositorio espejo local (Mirror) para contenido pesado (Wikipedia, Khan Academy, paquetes de software), evitando romper el cifrado punto a punto SSL/TLS innecesariamente para salvaguardar la privacidad de la telemedicina.
    \item \textbf{Directivas de Refresco (\texttt{refresh\_pattern}):} Se configuran tiempos mínimos de retención agresivos para contenido estático (PDFs, imágenes, scripts: \texttt{refresh\_pattern -i \textbackslash.(pdf|epub|mp4|jpg|css|js)\$ 10080 90\% 43200 override-expire}).
\end{itemize}

\subsection{Mecanismo de Conmutación Online/Offline y Sincronización}
Se implementa un script de detección de conectividad (Heartbeat a servidores DNS públicos 1.1.1.1 y 8.8.8.8):
\begin{itemize}
    \item \textbf{Modo Online Detectado:} Se activa la cola de sincronización en segundo plano (\texttt{rsync} diferido y actualización de repositorios críticos en horas de baja demanda).
    \item \textbf{Modo Offline Detectado:} Squid conmuta a modo autónomo (\texttt{offline\_mode on}), respondiendo exclusivamente con objetos persistidos en disco y evitando timeouts en peticiones DNS upstream.
\end{itemize}

\section{Cronograma de Proyecto y Ruta Crítica}

El cronograma general de ejecución se estructura en las tres fases del proyecto:

\begin{table}[htbp]
\caption{Cronograma de Fases y Actividades Principales}
\begin{center}
\begin{tabularx}{\columnwidth}{l l X}
\toprule
\textbf{Hito / Fase} & \textbf{Semanas} & \textbf{Entregables Clave} \\
\midrule
\textbf{Fase 1: Plan} & Sem. 1--3 & Documento de arquitectura IEEE, diagramas L2/L3, matriz de riesgos y especificación de hardware. \\
\textbf{Fase 2: Prototipo} & Sem. 4--7 & SO Linux afinado, contenedor Squid funcional, DHCP/DNS local, portal cautivo preliminar. \\
\textbf{Fase 3: Operación} & Sem. 8--11 & Servidor Jellyfin integrado, conmutación online/offline automática, dashboard categorizado y pruebas de carga. \\
\bottomrule
\end{tabularx}
\end{center}
\label{tab:schedule}
\end{table}

\textbf{Ruta Crítica:} La actividad crítica con mayor riesgo de cuello de botella es la configuración y afinamiento del almacenamiento masivo (2--4 TB) y la persistencia de caché HTTPS offline en Squid, ya que condiciona el éxito de la Fase 2 y 3.

\section{Gestión de Riesgos y Planes de Mitigación}

\begin{table}[htbp]
\caption{Matriz de Gestión de Riesgos Técnicos y Ambientales}
\begin{center}
\begin{tabularx}{\columnwidth}{l c X}
\toprule
\textbf{Riesgo Identificado} & \textbf{Impacto} & \textbf{Estrategia de Mitigación} \\
\midrule
Corte repentino de fluido eléctrico en la unidad móvil & Crítico & Incorporación de UPS de respaldo con señal de apagado ordenado (\texttt{nut/apcupsd}) y sistema de archivos journaling con montura segura (\texttt{data=ordered}). \\
Falta total de enlace satelital o móvil & Alto & Operación en modo aislado (\texttt{offline\_mode on}) con DNS local resolviendo hacia servidores internos (\texttt{dnsmasq} estático). \\
Corrupción o saturación de disco de 2--4 TB & Alto & Script periódico de mantenimiento de espacio con purga LRU (Least Recently Used) y partición separada para el sistema operativo. \\
Sobrecarga térmica del RPi5 o Mini PC & Medio & Disipador activo con ventilador PWM y monitoreo continuo de temperatura de CPU mediante script ligero. \\
\bottomrule
\end{tabularx}
\end{center}
\label{tab:risk_matrix}
\end{table}

\section{Conclusiones y Trabajo Futuro}
El plan de ingeniería propuesto para la Fase 1 establece las bases técnicas completas para responder a los criterios de evaluación de nivel superior (\textit{Accomplished}). Al combinar segmentación estricta L2/L3, políticas avanzadas de proxy/caché con Squid, tolerancia a fallos eléctricos y operación 100\% offline, se garantiza una solución viable, escalable y adaptada a las complejidades del entorno rural. En la Fase 2, se procederá a la construcción del prototipo automatizado mediante contenedores Docker.

\begin{thebibliography}{00}
\bibitem{squid} Squid Web Proxy Cache, ``Configuration Manual and Tuning Guidelines,'' \url{http://www.squid-cache.org/Doc/config/}, 2024.
\bibitem{opennds} openNDS Documentation, ``Next Generation Captive Portal Engine,'' \url{https://opennds.readthedocs.io/}, 2024.
\bibitem{rfc1918} Y. Rekhter et al., ``Address Allocation for Private Internets,'' BCP 5, RFC 1918, Feb. 1996.
\bibitem{ieee8021q} IEEE, ``IEEE Standard for Local and metropolitan area networks--Bridges and Bridged Networks,'' IEEE Std 802.1Q-2018.
\end{thebibliography}

\end{document}
